% ============================================================
% Section 127: 施主杂质电离与 Gunn 效应
% ============================================================
\section{半导体物理：施主杂质电离与 Gunn 效应}
\label{sec:donor-ionization-gunn}

\begin{modelbox}[题目：施主杂质能级与全电离的物理机制]
试画出掺有施主杂质 Ge、Si 等半导体能带图，为什么能这样画？施主杂质能级在导带底下面还是上面？如果在下面，杂质能级被电子占据几率较大，又如何解释杂质全电离时杂质原子上的电子几乎全呆在导带底的实验事实呢？

（补充讨论：应用类似的分析方法解释 GaAs 晶体中的电子转移效应——Gunn 效应。）
\end{modelbox}

\begin{tipbox}[考点快速分析]
\begin{itemize}
    \item \textbf{类氢模型}：施主杂质引入禁带中靠近导带底的分立能级
    \item \textbf{全电离的悖论}：能级低却被排空——态密度竞争导致``熵致电离''
    \item \textbf{Gunn 效应}：多谷能带结构下，强电场驱动电子从低态密度谷转移到高态密度谷
    \item \textbf{共同物理}：电子在不同态密度状态间的重新分布决定宏观性质
\end{itemize}
\end{tipbox}

\begin{derivationbox}[标准解题过程]
\textbf{1. 施主杂质能带图及其物理依据}

\textit{能带图}：在禁带中，靠近导带底 $E_c$ 下方的位置画出一条分立的短线，代表施主杂质能级 $E_D$。能级与导带底的能量差 $E_i=E_c-E_D$ 即为施主电离能（约 $0.01\sim0.04\,\text{eV}$）。

\textit{物理依据}：以硅、锗中掺磷（P）为例。磷为五价元素，取代晶格中的四价硅原子后，四个价电子与周围硅原子形成共价键，多余的一个价电子受到磷离子实（正电中心）的库仑吸引，形成类氢束缚态。该束缚态的能量低于导带底，因此在禁带中靠近导带底处引入一个分立的施主能级 $E_D$。

\textbf{2. 杂质能级被电子占据几率与全电离的矛盾及解释}

\textit{表面矛盾}：施主能级 $E_D$ 位于导带底 $E_c$ 下方。根据费米-狄拉克分布，能量越低的能级被电子占据的概率越大。若仅看单电子能级位置，似乎施主能级应被电子填满，而非电离。

\textit{全电离的实验事实}：在室温下，掺杂浓度适中的半导体中，几乎所有施主杂质都已电离，电子进入导带，施主能级几乎全空。

\textit{解释}：关键在于\textbf{态密度权重}的竞争。导带拥有巨大的有效态密度（$N_c\sim10^{19}\,\text{cm}^{-3}$），而施主能级的态密度仅为掺杂浓度 $N_D$（通常 $10^{14}\sim10^{17}\,\text{cm}^{-3}$）。在热平衡下，电子在施主能级和导带之间的分布由两者的态密度和能量差共同决定。

室温下 $k_BT\approx0.026\,\text{eV}$，而 $E_c-E_D\sim0.01\sim0.04\,\text{eV}$。由于 $N_c\gg N_D$，两者的乘积使得 $n/n_D\gg1$。因此，热平衡下绝大多数电子被``挤''到了态密度巨大的导带中，施主能级反而近乎全空。这就是\textbf{熵驱动}的电离：电子进入导带虽然能量略高，但可占据的微观状态数剧增，系统总自由能降低。

\textbf{3. 补充：GaAs 中的电子转移效应（Gunn 效应）}

\textit{GaAs 的能带结构}：
\begin{itemize}
    \item 导带底位于布里渊区中心 $\Gamma$ 点（中心能谷），有效质量小（$m_1^*\approx0.066m_0$），电子迁移率高。
    \item 在 $[100]$ 方向的 $L$ 点存在能量较高的卫星能谷，比中心能谷高约 $0.3\sim0.4\,\text{eV}$，有效质量大（$m_2^*\approx0.85m_0$），态密度大（约为中心谷的 46 倍）。
\end{itemize}

\textit{电子转移过程}：
\begin{itemize}
    \item 低电场时，电子聚集在中心能谷，迁移率高，电流随电场线性增长。
    \item 当电场足够强（$>3\,\text{kV/cm}$），电子从电场获得足够能量，通过谷间散射跃迁到卫星能谷。
    \item 卫星能谷的态密度远大于中心能谷，电子一旦进入，由于``态密度陷阱''效应，很难返回中心谷。大量电子聚集在卫星能谷，而这些电子的有效质量大、迁移率极低。
    \item 随着电场继续增大，越来越多的电子转移到低迁移率能谷，导致平均迁移率下降，电流反而随电场减小——出现\textbf{负微分电阻}。
\end{itemize}

负阻效应导致体内电荷畴的形成和渡越，产生微波振荡（Gunn 振荡）。

\textit{与杂质电离的相似性}：两者均为粒子在不同态密度权重的能级/能谷之间的重新分布。杂质电离是温度驱动的热平衡分布；Gunn 效应是电场驱动的非平衡输运过程，但``态密度陷阱''导致电子转移并滞留在高能谷的物理思想一致。
\end{derivationbox}

\begin{formulabox}[答案汇总]
\begin{center}
\begin{tabular}{ll}
\toprule
\textbf{问题} & \textbf{答案要点} \\
\midrule
施主能带图 & 禁带中靠近 $E_c$ 下方画分立的施主能级 $E_D$ \\
全电离解释 & 导带态密度远大于施主能级，电子为增加熵而跃入导带 \\
Gunn 效应机制 & 强电场下电子从低有效质量、低态密度的中心谷转移到 \\ & 高有效质量、高态密度的卫星谷，产生负微分电导 \\
\bottomrule
\end{tabular}
\end{center}
\end{formulabox}

\begin{insightbox}[物理图像]
\begin{itemize}
    \item 施主杂质全电离是统计力学中态密度竞争导致``熵致电离''的经典范例。能量最低原理被熵增原理压倒，因为导带提供了几乎无限的``座位''。
    \item GaAs 的 Gunn 效应是能带多谷结构导致电场驱动的电子转移和负阻振荡，是固态微波源（Gunn 二极管）的基础，工作频率可达 GHz 至 THz 量级。
    \item 两种现象均体现了固体中电子在不同态密度状态间的重新分布对宏观性质的深刻影响。
\end{itemize}
\end{insightbox}
